% !TeX root = ../papers/nc-p3.tex
\section{円錐曲線の極の平方座標と負付値の閉領域}
\label{sec:rank-fifteen-negative-x-fold}
\begin{theorem}\label{thm:rank-fifteen-negative-x-fold}
剰余標数3の付値で $10\nu(y)\le19\nu(x)<0$ となる標数零の通常な階数15窓は
標準Koszul左右AS正則であり、各始点で四中心四次式上自由階数256である。
\end{theorem}
\begin{proof}
$\nu(y)\le2\nu(x)$ は定理\ref{thm:rank-fifteen-negative-x-boundary}と既知の極領域による。
残る範囲では $s=i+w$, $r=gw$ の表示で $w,g$ はいずれも正付値である。
$B=2s/(1+s^2)$, $A=(1-s^2)/(1+s^2)$ として
\[
 h=1/B,\qquad z=h^2,\qquad u=rB,\qquad v=A/B
\]
と置く。$u/g$ と $h/w$ は単元であり、$v^2=z-1$, $v\equiv-i$ である。
倍率は $n=vu$, $H/n=u$, $K^2=1+zu^2$, $H=vu^2$ を満たす。
$AB=v/z$ より正則性座標は
\[
 x=(4+zu^2)(z-1)/z^2,\qquad y=u^4(z-1)^3/z^4
\]
となる。従って $h$ の符号を使わず、変数 $z,u$ だけで倍率を展開できる。

$K\equiv1$ とし、$P_\pm^2=2(K+1\pm vu^2)$ の根の定数項を2、
$(Q_\pm/u)^2=2z/(K+1)\pm2v$ の根の定数項を $1\mp i$ とする。
全て単純根である。$z=0$ の値と比べると $v$ と $Q_\pm/u$ の差は $(z)$ に、
$K$ と $P_\pm$ の差は $(zu^2)$ に入る。従って四次テンソルを与える各係数の差は
$(zu)$ に入る。これは整な局所環のイデアルの包含である。

\path{computations/verify_rank_fifteen_negative_x_fold.py} は $z=0$ で
$\mathbb Q(i)$ 上の $u$ の二十一桁を使う。未切断経路の損失は $(0,1,1,3,2)$、
四次中心は $(0,1,2,2)$、五次の積は二桁、六次から十三次は $(4,1,1,1,1,0,1,0)$
であり、還元可能な単元主元だけで最後の一桁が残る。
全308行列の指数と中心の自由語座標は前節の $w=g^2$ のものと一致する。
仮定は $\nu(w)\ge10\nu(g)$、従って $\nu(z)\ge20\nu(u)$ に等しい。
このときテンソルの差は $(u^{21})$ に入り、同じ桁数の消去をそのまま持ち上げられる。
経路、四次中心、十三次での中心商の消滅から既述の全次数の判定が従う。
\end{proof}
